\section{总结与延伸}

Hyung Won Chung 用 67 页 deck 完成了一次“架构考古”：先提出 exponentially cheaper compute 这一方向性力量，再用 Bitter Lesson 解释结构与规模的张力，随后把 encoder-decoder 到 decoder-only 拆成四个局部变换，最后用翻译、FLAN、层级表示与多轮缓存说明每项偏置为什么曾经合理、又为什么可能过期。

\subsection{十条核心结论}

\begin{enumerate}
  \item 研究快速变化领域时，追踪结果之外还要研究变化轨迹。
  \item Dominant driving force 是带误差边界的建模近似，不是唯一因果宣言。
  \item Compute per dollar 的历史趋势能指导研究策略，但不能自动证明未来斜率。
  \item Inductive bias 同时带来当前效率与潜在 scale ceiling。
  \item ``Scalable'' 指随预算扩大仍持续改善，不等同于能分布式运行。
  \item Encoder-decoder、encoder-only 与 decoder-only 共享 sequence-model 底座，差异可局部拆解。
  \item 四步 transformation 分别移除 attention 参数、stack 参数、layer connection 与 mask 的结构假设。
  \item 机器翻译和 long-input/short-target 数据解释了 encoder-decoder 为什么曾经有优势。
  \item Final-layer-only cross-attention 与 bidirectionality 的价值都应随深度和交互模式重新测量。
  \item 架构历史真正提供的是一套 bias audit，而不是永恒的赢家名单。
\end{enumerate}

\subsection{一张表压缩结构生命周期}

\begin{table}[H]
\centering
\small
\begin{tabular}{@{}L{0.17\textwidth}L{0.24\textwidth}L{0.24\textwidth}L{0.25\textwidth}@{}}
\toprule
结构假设 & 旧环境中的理由 & Regime 变化 & 应做的测试 \\
\midrule
输入/目标参数分离 & 翻译两侧角色和词表不同 & 通用知识、长生成、多轮对话 & 参数共享 scaling curve \\
只读最终 encoder layer & 中等深度下足够、实现简单 & 网络更深、粒度层级更丰富 & layer-aligned cross-attention \\
输入双向 attention & 完整输入上的理解任务 & streaming 与多轮增量交互 & 质量收益对重编码成本 \\
固定 architecture label & 便于实现与比较 & 组件可局部替换、系统约束变化 & 单变量 transformation ablation \\
\bottomrule
\end{tabular}
\caption{把历史结构写成可检验的生命周期假设。}
\end{table}

\subsection{实践：Bias audit 的最小流程}

下面的伪代码把课堂方法变成研究记录。它不会替代实验，却能防止团队只报告当前最好分数，而忘记记录结构成立的条件。

\begin{lstlisting}[caption={Inductive-bias audit 的研究记录模板}]
for bias in architecture.assumptions:
    document(original_problem(bias))
    measure(current_regime_gain(bias))
    choose_scale_axes(compute, data, depth, interaction)
    run_single_variable_ablation(bias, scale_axes)
    report(crossover_point, uncertainty, system_cost)
\end{lstlisting}

\begin{lstlisting}[caption={区分当前最优与长期趋势的决策伪代码}]
if deployment_budget < observed_crossover:
    ship(currently_efficient_method)
else:
    validate(less_structured_method)
keep_both_curves_and_failure_cases()
never_label_extrapolation_as_measurement()
\end{lstlisting}

\begin{warningbox}{不要把“减法研究”变成新的教条}
移除结构并不天然更科学。若某个偏置表达真实不变性、显著降低数据需求，或满足延迟、隐私与能耗约束，它可能长期合理。正确动作是测量收益、成本与交叉点，而不是把 ``less structure'' 当作无需证据的审美偏好。
\end{warningbox}

\subsection{自测问题}

\begin{enumerate}
  \item 为什么落笔实验允许忽略空气阻力，却不能推出复杂系统永远只有一个主导变量？
  \item Compute per dollar 与 Moore's law 有什么区别？课堂图没有证明什么？
  \item Inductive bias 如何同时提高 sample efficiency 又降低长期 scalability？
  \item Encoder-decoder 的三种 attention role 分别是什么？
  \item 四步 transformation 每一步移除了哪一项结构假设？
  \item 为什么机器翻译天然支持输入/目标参数分离？
  \item FLAN-T5 的更大增益为什么只是支持数据分布假设，而不是完成因果证明？
  \item Final-layer-only cross-attention 在极深网络中可能造成什么瓶颈？
  \item Bidirectional attention 为什么妨碍多轮对话中的 KV cache 复用？
  \item 如果 compute 趋势改变，应该更新哪些假设，而不是机械保留哪句结论？
\end{enumerate}

\subsection{拓展阅读}

\begin{itemize}
  \item Richard Sutton, \emph{The Bitter Lesson}：一般方法、计算与人工结构的长期关系。
  \item Vaswani et al., \emph{Attention Is All You Need}：原始 encoder-decoder Transformer 与三类 attention role。
  \item Devlin et al., \emph{BERT}：encoder-only 与 masked language modeling 的代表性设计。
  \item Brown et al., \emph{Language Models are Few-Shot Learners}：decoder-only scaling 的历史节点。
  \item Chung et al., \emph{Scaling Instruction-Finetuned Language Models}：FLAN-T5/PaLM 的 instruction-tuning 证据。
\end{itemize}

\end{document}
